\documentclass[12pt, a4paper]{article}

% Packages
\usepackage[utf8]{inputenc}
\usepackage[T1]{fontenc}
\usepackage[margin=2.5cm]{geometry}
\usepackage{amsmath, amssymb}
\usepackage{graphicx}
\usepackage{booktabs}
\usepackage{hyperref}
\usepackage{caption}
\usepackage{float}
\usepackage{listings}
\usepackage{xcolor}
\usepackage{parskip}
\usepackage{enumitem}
\usepackage{multirow}

% Figure search paths (works whether compiled from reports/ or repo root)
\graphicspath{
  {figures/}
  {../outputs/resnet18_potato_subset/}
  {./outputs/resnet18_potato_subset/}
}

% Hyperlink style
\hypersetup{
  colorlinks = true,
  linkcolor  = blue!60!black,
  urlcolor   = blue!60!black,
  citecolor  = blue!60!black,
}

% Code listing style
\lstset{
  basicstyle      = \ttfamily\small,
  breaklines      = true,
  frame           = single,
  rulecolor       = \color{gray!40},
  backgroundcolor = \color{gray!5},
  columns         = flexible,
}

% Title
\title{
  \Large\textbf{Knowledge Distillation for Edge-Friendly Plant Disease
  Classification}\\[0.4em]
  \large Progress Checkpoint Report\\[0.2em]
  \normalsize Machine Learning Term Project
}

\author{
  Meric Kelemeden\\
  \small\textit{Department of Computer Engineering}
}

\date{\today}

% ============================================================
\begin{document}
\maketitle
\tableofcontents
\newpage

% ============================================================
\section{Abstract}
% ============================================================

Deep convolutional neural networks achieve strong accuracy on image
classification tasks but are often too resource-intensive for edge deployment.
This project investigates knowledge distillation as a model-compression
technique for plant disease classification using the PlantVillage dataset.
A pretrained ResNet-18 is established as the initial supervised baseline,
achieving \textbf{98.8\%} test accuracy and a \textbf{97.9\%} macro-F1 score
on a three-class Potato disease subset after only three training epochs on CPU.
The distilled student model (MobileNetV3-Small) and the full 38-class
experiments are planned for the next phase.
All code, configuration files, and experiment scripts are maintained in the
project repository with full reproducibility.

% ============================================================
\section{Introduction}
% ============================================================

Automated detection of plant diseases from leaf images has the potential to
assist smallholder farmers who lack access to agronomists, particularly if
such systems can run on inexpensive mobile or edge hardware.
Deep learning models can match or exceed expert-level accuracy on curated
datasets~\cite{plantvillage2016}, yet models like ResNet-50 or EfficientNet
carry millions of parameters and impose latency costs that are impractical
on devices with limited compute and battery budgets.

\textbf{Knowledge distillation}~\cite{hinton2015distilling} addresses this
gap by training a compact \emph{student} network to mimic the output
distribution of a larger, more accurate \emph{teacher}.
Rather than learning only from hard one-hot labels, the student also learns
from the teacher's softened probability vectors, which encode inter-class
similarity and carry richer supervisory signal than a single correct-class
indicator.

This project asks:

\begin{quote}
  \emph{Can a distilled MobileNetV3-Small student for PlantVillage
  classification maintain competitive accuracy while improving parameter
  count, model size, and inference latency compared with its ResNet-50
  teacher?}
\end{quote}

The present report covers the \textbf{progress checkpoint}, which establishes
a reproducible data pipeline, a configurable model factory, and a supervised
ResNet-18 baseline on a fast-running Potato-disease subset.

% ============================================================
\section{Related Work}
% ============================================================

\subsection{Plant Disease Classification}

Mohanty et al.~\cite{plantvillage2016} demonstrated that a CNN trained on the
PlantVillage dataset could identify 26 diseases across 14 crop species with
over 99\% accuracy in controlled conditions.
Subsequent work has explored data augmentation~\cite{toda2019how},
semi-supervised learning, and transfer learning from ImageNet to improve
generalisation to real field conditions.

\subsection{Knowledge Distillation}

Hinton et al.~\cite{hinton2015distilling} introduced the concept of soft
targets with temperature scaling to transfer the generalisation capacity of a
large model to a smaller one.
The distillation objective blends a hard-label cross-entropy term with a
Kullback--Leibler divergence between softened student and teacher distributions:

\begin{equation}
  \mathcal{L} = \alpha \cdot T^{2} \cdot
    \mathrm{KL}\!\left(
      \sigma\!\left(\frac{z_s}{T}\right) \,\big\|\,
      \sigma\!\left(\frac{z_t}{T}\right)
    \right)
  + (1-\alpha)\cdot\mathrm{CE}(z_s, y)
  \label{eq:distill}
\end{equation}

where $z_s$ and $z_t$ are the student and teacher logits, $T$ is the
temperature, $\alpha$ is the soft-target weight, and $y$ is the ground-truth
label.
The $T^{2}$ factor compensates for the reduced gradient magnitude caused by
temperature scaling.

\subsection{Efficient Architectures for Edge Deployment}

MobileNets~\cite{howard2017mobilenets} use depth-wise separable convolutions
to drastically reduce multiply-accumulate operations.
MobileNetV3~\cite{howard2019searching} further incorporates hard-swish
activations and squeeze-and-excitation modules found by neural architecture
search, achieving a strong accuracy--latency trade-off on ImageNet.
EfficientNet~\cite{tan2019efficientnet} applies compound scaling to balance
depth, width, and resolution, offering another competitive teacher candidate.

% ============================================================
\section{Dataset and Preprocessing}
% ============================================================

\subsection{PlantVillage}

PlantVillage~\cite{plantvillage2016} is a publicly available dataset of
healthy and diseased leaf images across 38 plant--disease combinations.
Images are collected under controlled lighting against uniform backgrounds.
The full dataset contains approximately 54,000 RGB images; individual class
sizes vary considerably.

For the progress checkpoint, training is restricted to the three
\textbf{Potato} classes to enable rapid iteration:

\begin{itemize}[noitemsep]
  \item \texttt{Potato\_{}\_{}\_{}Early\_{}blight} \hfill (133 test samples)
  \item \texttt{Potato\_{}\_{}\_{}Late\_{}blight} \hfill (102 test samples)
  \item \texttt{Potato\_{}\_{}\_{}healthy} \hfill (23 test samples)
\end{itemize}

\subsection{Train / Validation / Test Split}

The dataset is not provided with pre-defined splits in the downloaded archive.
We apply a random split with a fixed seed (42) to obtain:

\begin{table}[H]
  \centering
  \caption{Dataset split sizes for the Potato subset (1,721 images total).}
  \label{tab:split}
  \begin{tabular}{lrr}
    \toprule
    \textbf{Split} & \textbf{Samples} & \textbf{Fraction} \\
    \midrule
    Train      & 1,205 & 70\% \\
    Validation &   258 & 15\% \\
    Test       &   258 & 15\% \\
    \bottomrule
  \end{tabular}
\end{table}

\subsection{Preprocessing and Augmentation}

All images are resized to $224 \times 224$ pixels and normalised using
ImageNet channel means and standard deviations
($\mu = [0.485,\,0.456,\,0.406]$,
$\sigma = [0.229,\,0.224,\,0.225]$).

Training images additionally receive:
\begin{itemize}[noitemsep]
  \item Random horizontal flip (probability 0.5)
  \item Random rotation up to $\pm 10^{\circ}$
\end{itemize}

Validation and test images use only resize and normalise transforms.
Labels are remapped to a contiguous $0 \ldots C-1$ range after subset
filtering so that the cross-entropy loss always receives valid targets.

% ============================================================
\section{Methodology}
% ============================================================

\subsection{Model Architecture}

The model factory (\texttt{src/models/\_\_init\_\_.py}) supports four
architectures loaded with pretrained ImageNet weights via the
\texttt{torchvision} weights API.
The final classification layer is replaced with a linear layer of size
$d_{\mathrm{in}} \times C$, where $C$ is the number of target classes.

\begin{table}[H]
  \centering
  \caption{Supported model architectures.}
  \label{tab:models}
  \begin{tabular}{llr}
    \toprule
    \textbf{Role} & \textbf{Model} & \textbf{Approx.\ Parameters} \\
    \midrule
    Baseline / Teacher candidate & ResNet-18       & 11.2 M \\
    Teacher                      & ResNet-50       & 25.6 M \\
    Teacher candidate            & EfficientNet-B0 &  5.3 M \\
    Student                      & MobileNetV3-Small & 2.5 M \\
    \bottomrule
  \end{tabular}
\end{table}

\subsection{Supervised Baseline Training}

The baseline training procedure (\texttt{src/train\_baseline.py}) follows a
standard supervised protocol:

\begin{itemize}[noitemsep]
  \item \textbf{Loss}: Cross-entropy
  \item \textbf{Optimiser}: Adam, $\mathrm{lr} = 10^{-3}$,
        weight decay $\lambda = 10^{-4}$
  \item \textbf{Epochs}: 3 (checkpoint scope)
  \item \textbf{Batch size}: 32
\end{itemize}

Device selection is automatic at runtime: CUDA $\succ$ Apple MPS $\succ$ CPU,
with no code changes required between machines.

\subsection{Distillation Loss (Next Phase)}

Equation~\ref{eq:distill} is already implemented in
\texttt{src/distillation.py}.
Planned hyperparameters are $T = 4.0$ and $\alpha = 0.5$.
Teacher and student will be trained on the full 38-class PlantVillage dataset
in subsequent experiments.

% ============================================================
\section{Experiments and Results}
% ============================================================

\subsection{Experiment Configuration}

All experiments are driven by YAML configuration files under
\texttt{configs/}.
The baseline is fully reproducible by running:

\begin{lstlisting}[language=bash]
python -m src.train_baseline --config configs/baseline_resnet18.yaml
\end{lstlisting}

Key hyperparameters for the current checkpoint are shown in
Table~\ref{tab:hparams}.

\begin{table}[H]
  \centering
  \caption{Hyperparameters for the ResNet-18 baseline experiment.}
  \label{tab:hparams}
  \begin{tabular}{lr}
    \toprule
    \textbf{Hyperparameter} & \textbf{Value} \\
    \midrule
    Model          & ResNet-18 (pretrained) \\
    Image size     & $224 \times 224$ \\
    Batch size     & 32 \\
    Epochs         & 3 \\
    Learning rate  & $1 \times 10^{-3}$ \\
    Weight decay   & $1 \times 10^{-4}$ \\
    Optimiser      & Adam \\
    Seed           & 42 \\
    \bottomrule
  \end{tabular}
\end{table}

\subsection{Training Dynamics}

Table~\ref{tab:history} shows per-epoch metrics on the training and
validation splits.
Loss converges rapidly; by epoch 3, validation accuracy and macro-F1 reach
100\% on the validation split of the three-class Potato subset, indicating
that pretrained ImageNet features transfer effectively to this domain.

\begin{table}[H]
  \centering
  \caption{Training history -- ResNet-18, Potato subset, 3 epochs.}
  \label{tab:history}
  \begin{tabular}{crrrr}
    \toprule
    \textbf{Epoch} &
    \textbf{Train Loss} &
    \textbf{Val Loss} &
    \textbf{Val Acc} &
    \textbf{Val Macro-F1} \\
    \midrule
    1 & 0.1634 & 0.3046 & 0.919 & 0.944 \\
    2 & 0.1122 & 0.2408 & 0.911 & 0.904 \\
    3 & 0.0459 & 0.0043 & 1.000 & 1.000 \\
    \bottomrule
  \end{tabular}
\end{table}

\subsection{Test Set Evaluation}

Table~\ref{tab:classreport} presents per-class precision, recall, and
F1-score on the held-out test split (258 samples).
The model misclassifies a single Early Blight sample and two Healthy samples,
the latter likely due to the small Healthy class support (23 test samples).

\begin{table}[H]
  \centering
  \caption{Per-class classification report on the test set (ResNet-18
  baseline).}
  \label{tab:classreport}
  \begin{tabular}{lrrrr}
    \toprule
    \textbf{Class} & \textbf{Precision} & \textbf{Recall} &
    \textbf{F1-score} & \textbf{Support} \\
    \midrule
    Potato Early Blight & 1.000 & 0.992 & 0.996 & 133 \\
    Potato Late Blight  & 0.971 & 1.000 & 0.986 & 102 \\
    Potato Healthy      & 1.000 & 0.913 & 0.955 &  23 \\
    \midrule
    Macro avg           & 0.990 & 0.969 & 0.979 & 258 \\
    Weighted avg        & 0.989 & 0.988 & 0.988 & 258 \\
    \bottomrule
  \end{tabular}
\end{table}

\subsection{Summary Metrics}

\begin{table}[H]
  \centering
  \caption{Baseline experiment summary -- ResNet-18, Potato subset.}
  \label{tab:summary}
  \begin{tabular}{lr}
    \toprule
    \textbf{Metric} & \textbf{Value} \\
    \midrule
    Test accuracy          & 98.84\%          \\
    Test macro-F1          & 97.88\%          \\
    Test weighted-F1       & 98.83\%          \\
    Trainable parameters   & 11,178,051       \\
    Inference latency      & 20.83 ms (CPU, batch=1) \\
    Training device        & CPU              \\
    Epochs                 & 3                \\
    \bottomrule
  \end{tabular}
\end{table}

\subsection{Confusion Matrix}

Figure~\ref{fig:cm} shows the confusion matrix on the test split.

\begin{figure}[H]
  \centering
  \includegraphics[width=0.55\textwidth]{confusion_matrix.png}
  \caption{Confusion matrix for ResNet-18 baseline on the Potato test set
  (258 samples, 3 classes).}
  \label{fig:cm}
\end{figure}

\subsection{Planned Next-Phase Comparison}

Table~\ref{tab:planned} will be completed in the next report once the
MobileNetV3 supervised baseline, the ResNet-50 teacher, and the distilled
student have been evaluated.

\begin{table}[H]
  \centering
  \caption{Planned experiment comparison table (dashes = not yet run).}
  \label{tab:planned}
  \begin{tabular}{llcccc}
    \toprule
    \textbf{Model} & \textbf{Training} &
    \textbf{Acc} & \textbf{Macro-F1} &
    \textbf{Params} & \textbf{Latency (ms)} \\
    \midrule
    ResNet-18           & Supervised   & 98.84\% & 97.88\% & 11.2M & 20.83 \\
    MobileNetV3-Small   & Supervised   & ---     & ---     & 2.5M  & ---   \\
    ResNet-50           & Supervised   & ---     & ---     & 25.6M & ---   \\
    MobileNetV3 (distilled) & Distillation & --- & ---  & 2.5M  & ---   \\
    \bottomrule
  \end{tabular}
\end{table}

% ============================================================
\section{Implementation Details}
% ============================================================

\subsection{Repository Structure}

\begin{lstlisting}
knowledge-distillation-edge/
+-- configs/
|   +-- baseline_resnet18.yaml
|   +-- student_mobilenetv3.yaml
|   +-- distillation.yaml        (next-phase placeholder)
+-- data/README.md               (download instructions)
+-- outputs/
|   +-- resnet18_potato_subset/  (metrics, CSV, confusion matrix)
+-- reports/report.tex
+-- scripts/
|   +-- run_baseline.sh
|   +-- run_student.sh
+-- src/
    +-- data/                    (dataset loading + label remapping)
    +-- models/                  (model factory)
    +-- train_baseline.py        (training entry point)
    +-- evaluate.py              (metrics + latency measurement)
    +-- distillation.py          (distillation loss, next phase)
    +-- utils/                   (config, seed, device helpers)
\end{lstlisting}

\subsection{Reproducibility}

\begin{itemize}[noitemsep]
  \item All experiments are seeded (Python, NumPy, PyTorch) via a single
        \texttt{seed} field in the YAML config (default: 42).
  \item No hard-coded file paths; \texttt{data\_dir} and
        \texttt{output\_dir} are config-driven.
  \item Automatic device selection runs identically on Apple Silicon (MPS)
        and NVIDIA CUDA without any code changes.
  \item Dataset download instructions are documented in
        \texttt{data/README.md}.
\end{itemize}

% ============================================================
\section{Conclusion and Future Work}
% ============================================================

This checkpoint establishes a clean, reproducible ML pipeline for plant
disease classification on the PlantVillage dataset.
The ResNet-18 supervised baseline achieves \textbf{98.84\% test accuracy}
and \textbf{97.88\% macro-F1} on the Potato subset after three training
epochs on CPU, confirming that ImageNet-pretrained features transfer
effectively to this domain.

The immediate next steps are:

\begin{enumerate}[noitemsep]
  \item Run the MobileNetV3-Small supervised baseline and record its
        accuracy, parameter count, and inference latency.
  \item Train a ResNet-50 teacher on the full 38-class PlantVillage dataset.
  \item Train a MobileNetV3-Small student via knowledge distillation using
        the loss in Equation~\ref{eq:distill} with $T=4.0$ and
        $\alpha=0.5$.
  \item Compare all models on accuracy, macro-F1, model size, parameter
        count, and single-image inference latency to quantify the
        compression--accuracy trade-off.
  \item Explore post-training quantisation (INT8) as an additional edge
        optimisation strategy.
\end{enumerate}

% ============================================================
\begin{thebibliography}{9}

\bibitem{plantvillage2016}
  S.~P. Mohanty, D.~P. Hughes, and M.~Salathe,
  ``Using Deep Learning for Image-Based Plant Disease Detection,''
  \textit{Frontiers in Plant Science}, vol.~7, p.~1419, 2016.

\bibitem{hinton2015distilling}
  G.~Hinton, O.~Vinyals, and J.~Dean,
  ``Distilling the Knowledge in a Neural Network,''
  \textit{NIPS Deep Learning Workshop}, 2015.

\bibitem{howard2017mobilenets}
  A.~G. Howard et al.,
  ``MobileNets: Efficient Convolutional Neural Networks for Mobile
  Vision Applications,''
  \textit{arXiv:1704.04861}, 2017.

\bibitem{howard2019searching}
  A.~Howard et al.,
  ``Searching for MobileNetV3,''
  \textit{IEEE/CVF ICCV}, pp.~1314--1324, 2019.

\bibitem{tan2019efficientnet}
  M.~Tan and Q.~V. Le,
  ``EfficientNet: Rethinking Model Scaling for Convolutional Neural
  Networks,''
  \textit{ICML}, 2019.

\bibitem{toda2019how}
  Y.~Toda and F.~Okura,
  ``How Convolutional Neural Networks Diagnose Plant Disease,''
  \textit{Plant Phenomics}, 2019.

\end{thebibliography}

\end{document}
